El cliente accede a la interfaz y puede hacer las siguientes opciones:
-	Crear nuevo lenguaje: puedes crear un lenguaje y dejarlo vacío
-	Editar lenguaje (add, delete, edit): para editar y borrar tiene que existir. No hay límite de palabras
-	Traducir (read): ningún limite. Si no entiende algo te da la opción de añadirlo a la BD


% ================================================================
%              FLUJO DE FUNCIONAMIENTO DEL SISTEMA
% ================================================================

COMPONENTES

1. Cliente: interactúa con la interfaz, graba los sonidos, crea lenguajes y los edita, recibe traducción
2. Coordinador: referido a todo lo relacionado con la interfaz y los otros componentes del sistema (API de espectogramas, filtros, BD...)
3. Audio Parser: API encargada de transformar sonidos en espectogramas. 
4. Base de datos: ¿qué tengo que explicar de la base de datos XD?
5. Intérprete: recibe una frases con las traducciones literales. Le pasa filtros y luego lo transforma en un lenguaje más natural.


CREAR LENGUAJE, EDITAR PALABRA, BORRAR PALABRA

Cliente ---> Interfaz ---> Coordinador ---> BD

En caso de error, la BD informaría al Coordinador


TRADUCIR

1. El cliente elige el lenguaje y le da a "traducir"
2. Indica que empieza una frase, reproduce toda la frase (sonidos) e indica que ha finalizado la frase.
3. Coordinador envía toda la frase al Audio Parser.
4. Audio Parser separa los sonidos en palabras. Devuelve un array de espectrogramas (respetando el orden de los sonidos) y null cuando detecta un sonido mayor de "t_min_silencio" segundos. Los demás silencios los elimina. 
5. Coordinador recibe el array de espectogramas y envía a la BD los espectrogramas, uno a uno. Sustituye los null por "---" o algo así, respetando en todo momento el orden de sonidos y silencios.
5.1. La BD devuelve una palabra (su significado). 
5.2. El Coordinador sigue enviando uno a uno los espectogramas hasta que el arary del espectogramas esté vacío. Si algún espectograma no tiene significado se considera como "unknonw" y se avisará al cliente. 
6. El Coordinador forma un array de significados (frase_traducida) y lo envía al Intérprete (sigue siendo Flutter)
7. El Intérprete filtra frase_traducida y forma la frase_traducida_filtrada.
8. Luego se comunica con un LLM (con un prompt establecido) y le envía frase_traducida_filtrada para que la vuelva más natural. Ejemplo: yo gustar gustar tú --> yo gustar (pasado) tú --> A mí me gustas tú. Crea frase_final
9. El LLM crea frase_final.
10. El Intérprete le devuelve la frase_final al Coordinador para que se lo reproduzca al cliente. 

AÑADIR PALABRA

1. El cliente edita el lenguaje y le da a "añadir nueva palabra".
2. Indica que empieza un sonido, lo reproduce y le da pausar. Indica el significado de ese sonido.
3. El Coordinador envía el sonido al Audio Parser.
4. El Audio Parser elimina los sonidos al inicio y al final (como un trim(), vaya) y devuelve el espectrograma del sonido.
5. El Coordinador lo recibe y le envía a la base de datos el espectrograma y el sonido.
6. La base de datos guarda la palabra
7. En caso de algún error, la BD avisa a Fluter y este se lo comunica al cliente (con una interfaz bonica, chicos)



% ================================================================
%                       ENTRADAS / SALIDAS 
% ================================================================

1. Clase epsectograma. Devolverá distancia coseno, MCFF y el otro metodo. La propia API podrá comparar aplicando distintas estrategias

2. Intérprete recibe un array de strings con los significados:
    [
    {
        "palabra": "gustar",
        "tipo": "verbo",
        "duracion": 0.5
    },
    {
        "palabra": "gustar",
        "tipo": "verbo",
        "duracion": 1.0
    }
    ]

    Donde pueden ser:
        - palabra: significado, "---" o ?
        - tipo: verbo, adverbio, adejtivo, pronombre, sustantivo, otros (silencios, signos, desconocido...)
        - duracion: proporción respecto a la duración guardada en la base de datos.

    En caso de que haya 1 silencio (t_min_silencio de duración)
    {
        "palabra": "---",
        "tipo": "otro",
        "duracion": 0
    }

    En caso de que no exista una traducción del sonido en la BD
    {
        "palabra": "unknown",
        "tipo": "otro",
        "duracion": 0
    }


% ================================================================
%                      REGLAS INTÉRPRETE
% ================================================================


1. Modificador por repetición inmediata

    1.1. Verbos: 
        - verbo --> verbo [presente o infinitivo]
        - verbo verbo --> verbo [pasado]
        - verbo verbo verbo --> verbo [futuro]
    1.2. Adverbio / Adjetivo:
        - adjetivo adjetivo --> adjetivo [mucho]
        - adjetivo adjetivo adjetivo --> adjetivo [extremo]
    1.3. Sustantivo:
        - sustantivo sustantivo --> sustantivo [plural]
    1.4. Pronombre: 
        - pronombre pronombre --> [posesivo: pronombre]
            Ejemplo: yo yo libro --> mi libro 

2. Puntuación

    2.1. "---" --> ","
    2.2. "---" "---" --> "."

3. Interrogación

    3.1. Frase ? --> ¿ Frase ?
    3.2. Frase ?? --> ¿ Frase ? [duda]
    3.3. Exclamación: Frase ! --> ¡ Frase !

4. Énfasis

    Si la duración de una palabra es > t_min_énfasis --> se pone en mayúsculas

5. Palabra desconocida

    Si recibe "unknown", literalmente introduce "palabra desconocida".

Limitar las repeticiones --> si una palabra se repite más de 3 veces, no se le aplicará ningún filtro. 



% ================================================================
%                        PALABRAS RESERVADAS
% ================================================================

- [] no se deben usar
- ---  no se debe usar 
- Si usamos "unknown" para sonidos sin significado conocido, el cliente no puede introducir un sonido que signifique "unknown" y que sea de tipo "otros"
- ¡ o ¿ no se deben usar


% ================================================================
%                          PRUEBAS
% ================================================================

- 4 pruebas de integración (1 para cada operación de CRUD) 
- Todas las pruebas componente que se puedan (zzzz)
- Todas las pruebas unitarias que se puedan (cuantas más mejor)



% ================================================================
%                      PATRONES APLICADOS
% ================================================================

-Patrón Estrategia en a API a la hora de decidir que algoritmo se aplica para
comparar los espectrogramas (cosineSimilarity, MCFF, ...)

